\label{sec:discrete-res}

\begin{figure}[htbp]
    \centering
    \includegraphics[width=\textwidth]{vocab-frequency-better.png}
    \caption{Normalised frequency of each code across 2880 runs.}
    \label{fig:vocab-frequency-better}
\end{figure}

The agent's utilisation of the vocabulary gives an entropy score of 4.61 and a Gini coefficient of 0.94 (Figure \ref{fig:vocab-frequency-better}).
This suggests the agent utilises the entire capacity of the codebook, unlike the AI Mother Tongue protocol.

\begin{figure}[htbp]
    \centering
    \includegraphics[width=\textwidth]{rank-frequency-vs-zipf-better.png}
    \caption{Vocabulary rank-frequency plot compared to Zipf's Law.}
    \label{fig:rank-frequency-zipf-better}
\end{figure}

The vocabulary usage closely follows Zipf's Law (Figure \ref{fig:rank-frequency-zipf-better}) until code around code 31, where, as before, the truncated Zipfian distribution is exhibited.

\begin{figure}[htbp]
    \centering
    \includegraphics[width=\textwidth]{comm_perplexity-better.png}
    \caption{Communication Perplexity tracked over a training run of WhisperWise trained across three runs. A 95\% confidence interval is shown.}
    \label{fig:comm-perplexity-over-training-better}
\end{figure}

The WhisperWise can learn to use the communication channel (Figure \ref{fig:comm-perplexity-over-training-better}), unlike the AI Mother Tongue protocol.
The protocol applies semantic compression during the first 5000 timesteps and further refines the language to a minimal set by about 20000 timesteps.
From then on, WhisperWise's communication protocol steadily increases, suggesting that the agents are learning to communicate the fine details of the environment.
Future work should evaluate how this communication protocol's perplexity varies over time until convergence.